\documentclass[11pt]{article}
\usepackage[T1]{fontenc}
\usepackage{lmodern,amsmath,amssymb,amsthm,booktabs,array}
\usepackage[margin=1in]{geometry}
\usepackage[colorlinks=true,linkcolor=blue,citecolor=blue,urlcolor=blue]{hyperref}
\newtheorem{theorem}{Theorem}
\newtheorem{corollary}[theorem]{Corollary}
\newcommand{\ac}{a_{\mathrm c}}
\newcommand{\ab}{\bar a}
\newcommand{\at}{a_{\mathrm T}}
\title{An alternative proof extending Kang's refutation of
Yang--Chu--Zhang (2015), Conjecture 1(ii), p.~15}
\hypersetup{pdftitle={An alternative proof extending Kang's refutation of Yang-Chu-Zhang (2015), Conjecture 1(ii), p. 15},pdfauthor={Haobo Ma and Wenlin Zhang}}
\author{Haobo Ma\\\small ChronoAI Pte Ltd, Singapore
\and Wenlin Zhang\\\small National University of Singapore\\
\small ChronoAI Pte Ltd, Singapore}
\date{November 7, 2026}
\begin{document}
\maketitle
\begin{abstract}
We give an alternative proof extending Kang's (2026) refutation of
Yang--Chu--Zhang (2015), Conjecture~1(ii), p.~15.
For its digamma remainder family $G_a$, we prove $G_a''(1/64)>0$
for every $a\geq331467377585/629168992256$, and $G_a''(1/64)>1$
for every $a\geq5269/10000$.
These bounds exclude complete monotonicity of $-G_a$ on a parameter
half-line containing strictly interior points of the conjectured range.
We also give a necessary positive constant-coefficient condition in the
authors' Problem~1 and characterize all its degree-four solutions as
positive multiples of the function in their Theorem~4 at its endpoint.
\end{abstract}

\section{The conjecture and the quantitative extension}
Yang, Chu and Zhang~\cite[Conjecture~1(ii), p.~15]{ycz} conjectured an
exact complete-monotonicity threshold. Kang~\cite{kang} refuted its
sufficient direction; we extend the explicit derivative obstruction and
deduce coefficient consequences for the accompanying Problem~1.
Write $\psi=\Gamma'/\Gamma$ and, for $x>0$, set
\begin{align}
R(x)&=\psi(x+1/2)-\log x,\label{eq:R}\\
P_a(x)&=x^4+ax^2+\frac{7a}{40}-\frac{31}{336},\nonumber\\
G_a(x)&=24P_a(x)R(x)-\left(x^2-\frac7{40}+a\right).
\label{eq:G}
\end{align}
These are the source's equations (1.10) and (3.3). Its conjecture asserts
\begin{equation}
-G_a\text{ is completely monotone on }(0,\infty)
\quad\Longleftrightarrow\quad a\leq\frac{155}{294}.
\label{eq:claim}
\end{equation}
The earlier posing is Yang~\cite[Conjecture~1(ii), p.~14]{yang}.
Complete monotonicity means smoothness and
$(-1)^nf^{(n)}(x)\geq0$ for every integer $n\geq0$ and every $x>0$.

Kang's Theorem~1 gives a strictly interior positive parameter and a
midpoint-convexity obstruction. His accompanying \texttt{PROOF.md}
also gives a derivative counterexample at $a=527/1000$, $x=1/50$.
The following estimate supplies a different evaluation point and an
explicit lower exclusion bound; it is not an exact replacement threshold.

\begin{theorem}\label{thm:main}
Put
\[
\ac=\frac{331467377585}{629168992256},\qquad
M=\frac{21943673}{1280},\qquad \ab=\frac{5269}{10000}.
\]
For every real $a\geq\ac$,
\begin{equation}
G_a''(1/64)>M(a-\ac)\geq0.\label{eq:main}
\end{equation}
For every $a\geq\ab$ the stronger numerical margin is
\begin{equation}
G_a''(1/64)>\frac{26102771679}{22937600000}>1.
\label{eq:simple}
\end{equation}
Thus $-G_a$ is not completely monotone whenever $a\geq\ac$.
In particular,
$0<\ac<\ab<527/1000<155/294$, so \eqref{eq:claim} is false.
\end{theorem}

\section{A rational second-derivative estimate}
The classical identities used below are~\cite[\S\S5.4(ii), 5.15]{dlmf}
\begin{equation}
\psi(1/2)=-\gamma-2\log2,\quad
\psi'(z)=\sum_{n=0}^{\infty}\frac1{(n+z)^2},\quad
\psi''(z)=-2\sum_{n=0}^{\infty}\frac1{(n+z)^3}\quad(z>0).
\label{eq:psi}
\end{equation}
The derivative series converge uniformly on positive compact intervals,
which justifies the last identity. Let $x_0=1/64$ and $z_0=33/64$.
We claim
\begin{equation}
R(x_0)>\frac{11}{6},\qquad R'(x_0)\geq-64,\qquad
R''(x_0)>4078.\label{eq:three}
\end{equation}
Indeed,
\[
\log2=2\int_0^{1/3}\frac{dt}{1-t^2}>\frac23,
\qquad \gamma<H_2-\log2<\frac56.
\]
Here $H_n-\log n$ decreases to $\gamma$, since
$\log(1+1/n)=\int_n^{n+1}dt/t>1/(n+1)$.
Positivity of $\psi'$ and $z_0>1/2$ give
$R(x_0)>4\log2-\gamma>11/6$, and
$R'(x_0)=\psi'(z_0)-64\geq-64$.
For the cubic series, the identity
\[
(n+97/64)^3-(n+1)(n+2)
=n^3+\frac{227}{64}n^2+\frac{15939}{4096}n
 +\frac{388385}{262144}>0\quad(n\geq0)
\]
implies
\[
\sum_{n=0}^{\infty}\frac1{(n+z_0)^3}
\leq\left(\frac{64}{33}\right)^3+
\sum_{n=0}^{\infty}\frac1{(n+1)(n+2)}
=\frac{298081}{35937}<9.
\]
The comparison series telescopes. Hence
$R''(x_0)=4096+\psi''(z_0)>4096-18=4078$.

\begin{proof}[Proof of Theorem~\ref{thm:main}]
Twice differentiating \eqref{eq:G} gives
\begin{equation}
G_a''(x)=24\bigl(P_a''(x)R(x)+2P_a'(x)R'(x)
+P_a(x)R''(x)\bigr)-2,\label{eq:derivative}
\end{equation}
where $P_a'=4x^3+2ax$ and $P_a''=12x^2+2a$.
At $a=\ac$, exact substitution gives
\[
P_{\ac}(x_0)=\frac{236541533}{3775013953536}>0.
\]
Since $\ac>0$, both $P_{\ac}'(x_0)$ and $P_{\ac}''(x_0)$ are
positive. All three quantities increase with $a$, with respective
slopes $3589/20480$, $1/32$ and $2$.
Thus all coefficients multiplying the bounds in \eqref{eq:three}
are nonnegative for $a\geq\ac$. Consequently
\[
G_a''(x_0)>L(a):=
24\left(\frac{11}{6}P_a''(x_0)-128P_a'(x_0)
+4078P_a(x_0)\right)-2.
\]
Rational arithmetic gives $L'(a)=M$ and $L(\ac)=0$, proving
\eqref{eq:main}, including the strict inequality at $a=\ac$.
Likewise $L(\ab)=26102771679/22937600000>1$, proving
\eqref{eq:simple}. Complete monotonicity of $-G_a$ would require
$(-G_a)''(x_0)\geq0$, a contradiction.
\end{proof}

\section{Coefficient consequences for Problem 1}
Yang--Chu--Zhang~\cite[Problem~1, p.~16]{ycz} ask for coefficients such that
\begin{equation}
F(x)=R(x)\sum_{k=0}^{n+1}a_kx^{2k}-\sum_{k=0}^{n}b_kx^{2k}
\quad\text{is completely monotone},\qquad
\lim_{x\to\infty}x^{2n+4}F(x)=c\ne0,\pm\infty.
\label{eq:problem}
\end{equation}
The finite nonzero limit is part of the question, in addition to
complete monotonicity.

\begin{corollary}\label{cor:constant}
Every solution of \eqref{eq:problem} has $a_0>0$.
More generally, a nonconstant completely monotone function of the
displayed polynomial--$R$ form has $a_0>0$.
\end{corollary}
\begin{proof}
Write $A(x)=\sum a_kx^{2k}$ and $B(x)=\sum b_kx^{2k}$.
The function $h(x)=\psi(x+1/2)$ is smooth near zero, so
$R(x)=-\log x+O(1)$ and $R'(x)=-1/x+O(1)$ as $x\downarrow0$.
If $a_0<0$, then $F(x)\to-\infty$, contradicting $F\geq0$.
If $a_0=0$, evenness gives $A(x)=O(x^2)$, $A'(x)=O(x)$ and
$B'(x)=O(x)$, whence $F'(x)=O(x|\log x|)+O(x)\to0$.
Complete monotonicity gives $F'\leq0$ and $F''\geq0$.
As $F'$ is increasing, its zero limit at the left endpoint forces
$F'\geq0$ as well. Thus $F$ is constant.
In \eqref{eq:problem}, a constant must be zero to have a finite scaled
limit, which would give $c=0$. This is excluded.
\end{proof}

\begin{corollary}\label{cor:quartic}
In Problem~1 with $n=1$, all solutions, and no others, are
\begin{equation}
F(x)=tG_{\at}(x),\qquad t>0,\qquad
\at=\frac{11165}{8284}.
\label{eq:quartic}
\end{equation}
Equivalently, their coefficients are
\[
(a_0,a_1,a_2)=24t\left(\frac{7\at}{40}-\frac{31}{336},\at,1\right),
\qquad (b_0,b_1)=t\left(\at-\frac7{40},1\right).
\]
Their scaled limit is
$c=627404761t/10287535104>0$.
\end{corollary}
\begin{proof}
The digamma asymptotic expansion~\cite[5.11.2]{dlmf}, translated by
$1/2$, gives
\begin{equation}
R(x)=\frac1{24x^2}-\frac7{960x^4}+\frac{31}{8064x^6}
-\frac{127}{30720x^8}+\frac{511}{67584x^{10}}+O(x^{-12}).
\label{eq:asymptotic}
\end{equation}
For $F=R(a_2x^4+a_1x^2+a_0)-(b_1x^2+b_0)$,
cancellation of the $x^2,x^0,x^{-2},x^{-4}$ coefficients is necessary
and sufficient for $F=O(x^{-6})$. These four equations give, successively,
\[
b_1=\frac{a_2}{24},\quad b_0=\frac{a_1}{24}-\frac{7a_2}{960},\quad
a_0=\frac{7a_1}{40}-\frac{31a_2}{336},\quad
a_1=\frac{11165}{8284}a_2.
\]
Thus $F=(a_2/24)G_{\at}$. The next coefficient in
\eqref{eq:asymptotic} gives the stated $c$ with $t=a_2/24$.
Since $F\geq0$ and $c\ne0$, necessarily $c>0$, hence $t>0$.
Conversely, Yang--Chu--Zhang's Theorem~4, pp.~12--13, proves that
$G_a$ is strictly completely monotone exactly when $a\geq\at$.
At $a=\at$ it supplies complete monotonicity, and the expansion supplies
the required finite positive limit. This proves sufficiency.
\end{proof}

The classification is a coefficient consequence of the source's
Theorem~4, not a new proof of that theorem. If the indexing allows $n=0$,
the same cancellation argument gives exactly
$F(x)=t\{24(x^2+7/40)R(x)-1\}$, $t>0$, with
$c=2071t/33600$; complete monotonicity is the source's Theorem~2.
Indeed $b_0=a_1/24$, $a_0=7a_1/40$, and $t=a_1/24$.

For comparison, substitution of the negatives of the coefficients in
\eqref{eq:G} gives $-G_a$, which Theorem~\ref{thm:main} excludes for
$a\geq\ac$. But most of that family does not satisfy Problem~1's
decay condition even before complete monotonicity is considered:
\begin{equation}
\lim_{x\to\infty}x^4(-G_a(x))
=-\frac{2071}{33600}(a-\at).
\label{eq:decay}
\end{equation}
Only $a=\at$ can have the required finite $x^6$ limit; there
$-G_{\at}$ has a negative limit coefficient and is excluded.
Thus the derivative counterexample must not be confused with a full
classification of the coefficient problem.

\section{Coverage and mathematical evidence}
\noindent
\begin{tabular}{@{}p{.30\linewidth}p{.65\linewidth}@{}}
\toprule
Source item in \cite{ycz} & Coverage\\
\midrule
Conjecture 1(ii), p.~15 & Refuted: Kang (2026). Theorem~\ref{thm:main}
extends the derivative exclusion. A corrected necessary condition is
$a<\ac$, not a replacement equivalence.\\[3pt]
Conjecture 1(i), p.~15 & Ratio monotonicity: not covered.\\[3pt]
Remark 3, pp.~14--15 & Its necessary condition $a\leq155/294$ is
strengthened to $a<\ac$. Bounds (4.2)--(4.3) do not imply complete
monotonicity.\\[3pt]
Problem 1, p.~16 & Answered for $n=1$ by
Corollary~\ref{cor:quartic}, using the prior Theorem~4;
Corollary~\ref{cor:constant} gives $a_0>0$ for every $n$.
The classification for $n\geq2$ is not covered.\\[3pt]
Theorems 2--4; Remarks 4--5 & Prior results cited, not refuted.
The endpoint in Theorem~4 supplies the positive family above;
the sequence bounds in Remarks 4--5 are unchanged.\\
\bottomrule
\end{tabular}

\medskip
The introductory discussion of Chen's conjectures on p.~3 concerns
different shifted functions; the source already refutes the Bernstein
claim for $H$ there and proves its cited complete-monotonicity claim
in Theorem~1. The other Bernstein claim and harmonic-number questions
are not covered by the present estimates.
The complete parameter set for $-G_a$ below $\ac$ is not determined here.

\paragraph{Written proofs, finite calculations and Lean declarations.}
In
\href{https://github.com/the-omega-institute/trureturing/blob/074df09ddbbef1843ddb383577bfaa827f753e36/D5/S3/Analytic/YangChuZhangPsiCM/Refutation.lean}{\texttt{Refutation.lean}}
at revision \texttt{074df09}, the declaration \texttt{result} proves the
negation of \eqref{eq:claim}; its definition \texttt{claim} quantifies
over every real $a$. The file's private lemma
\texttt{uniform\_\allowbreak second\_\allowbreak deriv} proves
$G_a''(1/64)>1/2$ for $a\geq155/294-1/10000$.
The lower exclusion bound $\ac$, the margin at $\ab$ and the coefficient
corollaries have written proofs here; they are not conclusions of that
Lean lemma. The standard-library program \texttt{verify.py} checks
the rational bounds, affine identities, translated asymptotic coefficients
and coefficient cancellations. Analytic identities, convergence,
endpoint reasoning and complete monotonicity supplied by the cited
Theorems~2 and~4 have the written or cited proofs specified above.

\begin{thebibliography}{9}
\bibitem{ycz}
Z.-H. Yang, Y.-M. Chu and X.-H. Zhang,
Necessary and sufficient conditions for functions involving the psi function
to be completely monotonic,
\emph{Journal of Inequalities and Applications} \textbf{2015}, 157 (2015).
\href{https://doi.org/10.1186/s13660-015-0674-8}{doi:10.1186/s13660-015-0674-8}.

\bibitem{yang}
Z.-H. Yang, Several completely monotone functions related to DeTemple's
sequence, 2014. \href{https://arxiv.org/abs/1409.6411v1}{arXiv:1409.6411v1}.

\bibitem{kang}
M. Kang, A failure of the DeTemple complete-monotonicity threshold,
2026. Theorem~1 and accompanying \texttt{PROOF.md}.
\href{https://doi.org/10.5281/zenodo.23718266}{doi:10.5281/zenodo.23718266}.

\bibitem{dlmf}
NIST Digital Library of Mathematical Functions,
\href{https://dlmf.nist.gov/5.4.E13}{equation 5.4.13},
\href{https://dlmf.nist.gov/5.15.E1}{equation 5.15.1} and
\href{https://dlmf.nist.gov/5.11.E2}{equation 5.11.2}.
Accessed November 7, 2026.
\end{thebibliography}

\smallskip\noindent\small
\copyright\ 2026 Haobo Ma and Wenlin Zhang.
Licensed under \href{https://creativecommons.org/licenses/by/4.0/}{CC BY 4.0}.
\end{document}
